% Transcription — fonds Alexandre Grothendieck, shelfmark 18, batch 2 (pages 21-40).
% Established from the facsimile archives/batches/18/batch-02.pdf (a mirror of
% https://grothendieck.umontpellier.fr/18.pdf), per the skill
% transcribe-grothendieck. The batch file carries no cover sheet: its PDF page 1
% is archivists' page 21. This was checked against the pencilled numbers at the
% bottom-left of the leaves — 24 on the fourth, 28 on the eighth, then 31, 32,
% 33, 34, 35, 36, 37, 38, 39 and 40 on the last ten, so the alignment holds end
% to end.
%
% Pass: Opus 5 (claude-opus-5), 2026-09-16 — first pass, unchecked against the pages by a human.
%
% Fourteen leaves carry his hand: 22, 24, 26, 28, then 30 (a cover) and 31-40
% unbroken. The five odd leaves 21, 23, 25, 27, 29 are the backs of typescripts
% foreign to these notes — a French typescript on formal kernels on 21, a
% Bourbaki « Tribu » n° 59 on root systems on 23, 25, 27 and 29, several of them
% scanned upside down and none of them marked by him. They are skipped and the
% gaps in the numbering are the record. Leaves 31 and 34 are written on Institut
% des Hautes Études Scientifiques letterhead; nothing on any leaf carries a date,
% and there is no pagination of his.
%
% What the batch is. Two runs.
%
% The first, pages 22-28, finishes the comparison of analytic and topological
% fundamental classes that occupies pages 14-20 of batch 1, then turns it into a
% question about the base field.
%   22  the fundamental formula of pages 16-20 generalised to an arbitrary proper
%       morphism f : X → Y of relative codimension p, as a boxed square comparing
%       f_*^top on H^*( ,C) with (1/(2iπ)^p) f_*^an on H^*( ,Ω^·).
%   24  the same compatibility for a correspondence ξ ∈ a^n(X×Y): the map
%       x ↦ pr_{2*}(pr_1^*(x)·P_top(ξ)) against the algebraic one, through
%       P_top(ξ) = (1/2iπ)^n P_alg(ξ); the canonical isomorphism
%       ξ_C(H^i(X,Q(0))) = ξ(H^i(X,Q(0))) ⊗_Z C ⥲ H^i(X,Ω^·_X); closing « Ouf ! »
%   26  a topological field K isomorphic to C but with NO isomorphism chosen, and
%       the functor ξ = ξ_K on motives over K with values in Z-modules, together
%       with ξ ⊗_Z Z_ℓ ⥲ ξ_ℓ and ξ ⊗_Z K ⥲ ξ_K; then ξ(H^2(P^1,Z(0))) = ξ(Z(-1)),
%       isomorphic to Z but not canonically — the choice amounts to an
%       orientation of P^1_K, i.e. to a square root of -1 in K.
%   28  the choice made, α_∞(u) = (1/(2√(-1)π))·v against the canonical generator
%       of H^2(P^1,Ω^1); and for K isomorphic to R (where the isomorphism IS
%       unique) a functor to Z-modules carrying an action of
%       Z/2Z = Gal(R-bar/R), inside ∏_ℓ M_ℓ × M_∞C.
%
% The second, pages 30-40 under his own cover, is real Hodge theory: what a
% topological base field K ≃ R or C does to the realisations of a motive, and
% what survives when no isomorphism K ≃ C is chosen.
%   30  cover: « Corps de base typ. ≃ R ou C (foncteurs de Hodge et De Rham dans
%       ce cas) ».
%   31  X smooth projective over K, X(K) a compact K-analytic manifold, hence the
%       Betti functor; Betti(M) ⊗_Q K ≃ H^*(X,K) and, by the usual argument,
%       Betti(M)_K ≃ DR(M) canonically as ⊗-isomorphism. This puts a real
%       structure on DR(M) whose real points are Betti(M)_R; the filtration and σ
%       together split it into a bigrading compatible with total degree.
%   32  the bigrading obtained by transport of structure is NOT independent of
%       the chosen isomorphism u : K ≃ C (u against ū); for K_0 = R, DR(M)_K
%       carries two R-structures, σ from the complex conjugation and τ from
%       DR(M)_K ≃ DR(M_K), and σ and τ commute, τ being the automorphism of
%       H^*(X(K),R) induced by the conjugation τ_X of X(K).
%   33  what a motive over R therefore gives: a Hodge structure V and a « real »
%       antilinear automorphism of V_K exchanging the filtration with its
%       opposite — the \emph{Frobenius réel}, an element f_R ∈ G(Q) of order 2
%       whose image generates G/G^0 and which normalises i : S_R → G_R with
%       f_R i(λ) f_R^{-1} = i(λ'), so f_R i(λ) f_R^{-1} = i(λ^{-1}) on U_R; the
%       functor M ↦ (Hodge structure, f_R) is fully faithful.
%   34  the elements z = x + iy of DR(M)_K ≃ V ⊗_Q K with x = fx, y = -fy, i.e.
%       τz = z; the space W they form, the « Betti tordu » of DR(M), and the
%       square H^*(X(K),Q) = V, H^*(X(R),Q) = V ∩ W = V^f; a second fibre functor
%       over Q, isomorphic to the first, the isomorphism V ≃ W being the identity
%       on V^+_K and multiplication by i on V^-_K — the two differ by the cocycle
%       of f_R.
%   35  (blue ink) ᾱ = f_∞ α, so f_∞ = 1 ⟺ α ∈ P(R) ⟺ R(α) = R; the boxed
%       ε(f_∞)ε(f_∞) = -1, hence f_∞ ≠ 1 and R(α) = C as soon as M contains the
%       Tate motive; and the equivalence (f_∞ = 1) ⟺ (R(α) = R) ⟺ M generated by
%       the Q(-2n) ⟺ G ≃ G_m or G ≃ {e}.
%   36  « Motifs sur R »: G = Aut_⊗(T_B), the ⊗-automorphism of order 2 of T_B
%       induced by x ↦ x̄ on X(C), giving f_∞ ∈ G(Q) with f_∞² = 1; the relation
%       τ_M = (f_∞)_M σ_M = σ_M (f_∞)_M; the torsor
%       P = Isom_⊗(T_B ⊗_Q R, T_DR) over G_R as the one defined by the 1-cocycle
%       on Γ = Z/2Z with value f_∞, the twisted form of G being int(f_∞) ×_R σ_C.
%   37  R(γ) = R(β,γ) = R(β,α) = R(α), which is C outside the exceptional case;
%       a Remarque that f_{∞M} = id_M forces M to be a sum of tensor powers of
%       Q(2); then a long computation in Norm_{G_R}(C,C^{-1}) cancelled by
%       diagonal strokes.
%   38  for K topological ≃ C and X over K, the isomorphism
%       H^*(X_top,K) ⥲ H^*(X) does not depend on the chosen u : K ≃ C; the two
%       homomorphisms for K = C differ by conjugation of the coefficients; hence
%       H^*(X_top,R) ↪ H^*(X), and for X projective non-singular a decomposition
%       φ : ⊔_{p,q} H^q(X,Ω^p_X) ⥲ H^*(X) compatible with the filtration and
%       independent of u.
%   39  the same over K ≃ R, with K' an algebraic closure and π = Z/2Z: the
%       isomorphisms commute with π, so descend; in short, for X projective
%       non-singular over a locally compact archimedean field one has
%       H^{p,q}(X) = φ(H^q(X,Ω^p_X)) ⊂ F^p H^q(X).
%   40  for K not topological, each archimedean topology T gives
%       H^*(X)_T = H^*(X) ⊗_K K̂_T and a boxed φ_T, whence the K-subspaces
%       H^{p,q}_T(X) and H^{p,q}_T(X,R); R_T = R ∩ K when K̂_T ≃ C, and
%       K_r ≃ ∩_T R_T, the largest totally real subfield of K.
%
% Notation fixed for the batch, kept in line with batches 1 and 3 of the folder.
% Z, Q, R, C, Z_ℓ, Q(0), Q(-n) are set \mathbb{}; K, K', K_0, K_r, X, Y, M, V, W
% stay roman. The fibre functors are T_B and T_DR as batch 1 writes them, and the
% functor of pages 26-28 is his ξ, which he writes as a tall script zeta —
% \xi throughout, with ξ_ℓ, ξ_K, ξ_{C,∞}. The two realisation functors of the
% second run are his own words Betti( ) and DR( ), set \mathrm{}; his « H^* » is
% written without the hypercohomology underlining batch 1 records, and the De
% Rham complex is \Omega^{\bullet}_{X} or, where he writes it with a star,
% \Omega^{*}_{X}. Cycle-class and Poincaré maps keep batch 1's labels P_top,
% P_alg, P_an, and a^n(X×Y) is the group of cycle classes, drawn with his rounded
% a. The Frobenius at the real place is f_R on pages 33 and 36 and f_∞ on pages
% 35-37; both are transcribed as he writes them and are the same element, which
% page 36 says in as many words. S_R is Serre's torus and U_R its circle
% subgroup; i and j are the two letters he uses for the same map S_R → G_R
% (page 33 writes i, page 36 writes j), and both are kept. G^0 is the identity
% component, G^δ appears once on page 35.
%
% His « ré » is drawn so that « réseaux » reads « viscaux » and « réel » can read
% « viel »; the reading is fixed once and held. « aut. » is automorphisme, and
% his half-formed « autom » is expanded with \add. The blue-ink leaves 35-37 are
% written fair and read almost end to end; 24, 26 and 28 are a fast palimpsest
% whose formulas carry and whose connective prose largely does not, and they
% carry most of this batch's apparatus.
%
% Batch boundaries. A run crosses the 20/21 boundary in substance though not in
% syntax: page 20 (batch 1) closes on a boxed square and page 21 is typescript,
% but page 22 opens « On généralise la formule fondamentale », and what it
% generalises is the formula of pages 16-20. A \note marks this at page 22.
% Nothing crosses 40/41: page 40 closes on an NB about K_r, and page 41 (batch 3)
% is a fresh cover in his hand naming a different run.

\documentclass[11pt,a4paper]{article}
% Le préambule commun à toutes les transcriptions du fonds Grothendieck.
%
% Il tient en une idée : **l'incertitude se compose, elle ne se résout pas.**
% Une transcription qui lisse un mot douteux a détruit la seule chose pour
% laquelle on la faisait — un lecteur doit pouvoir savoir, à chaque endroit,
% ce qui était sur le feuillet et ce qui a été ajouté. Les six macros qui
% suivent sont donc l'appareil critique, et non de la décoration.
%
% Les mêmes commandes sont reconnues par `scripts/render.mjs`, qui produit la
% vue de lecture du site. Une macro ajoutée ici doit l'être aussi là-bas,
% faute de quoi le rendu échouera bruyamment — ce qui est voulu : un rendu
% silencieusement faux est pire qu'un rendu absent.

\ProvidesPackage{grothendieck}[2026/08/08 Transcriptions du fonds Grothendieck]

\usepackage{amsmath,amssymb,amsthm}
\usepackage{mathrsfs}
\usepackage{stmaryrd}
% Les diagrammes commutatifs sont partout dans ce fonds : c'est la langue
% même de la géométrie algébrique de Grothendieck, et une transcription qui
% les rendrait en prose ne transcrirait rien.
\usepackage{tikz-cd}
% Français par défaut — c'est la langue des feuillets — mais l'anglais est
% chargé aussi : la traduction et le résumé sont des documents anglais, et la
% typographie française (espaces avant les deux-points) y serait une faute.
% Ces documents basculent par \selectlanguage{english} en tête de corps.
\usepackage[english,french]{babel}
\usepackage{fontspec}
\usepackage{geometry}
\geometry{a4paper,margin=25mm,marginparwidth=28mm}
\usepackage{marginnote}
\usepackage{xcolor}
% ulem, not soul. `soul`'s \st and \ul are built on a character-by-character
% scan that XeLaTeX's font machinery breaks: the document fails with an
% undefined control sequence rather than degrading. ulem does the same two jobs
% and is Unicode-engine safe. `normalem` stops it hijacking \emph, which the
% transcriptions use for Grothendieck's own emphasis.
\usepackage[normalem]{ulem}
\usepackage{hyperref}
\hypersetup{colorlinks=true,linkcolor=black,urlcolor=black,pdfborder={0 0 0}}

% --- Métadonnées du lot -----------------------------------------------------
% Reprises telles quelles par le script de rendu, qui les affiche en tête de
% la vue de lecture. Elles fixent ce que le fichier prétend couvrir : sans
% elles, une transcription flotte sans point d'attache dans un fonds de seize
% mille feuillets.

\newcommand{\folder}[1]{\gdef\grothFolder{#1}}
\newcommand{\batch}[1]{\gdef\grothBatch{#1}}
\newcommand{\leaves}[2]{\gdef\grothFirst{#1}\gdef\grothLast{#2}}
\newcommand{\foldertitle}[1]{\gdef\grothFoldertitle{#1}}
\gdef\grothFolder{?}\gdef\grothBatch{?}\gdef\grothFirst{?}\gdef\grothLast{?}\gdef\grothFoldertitle{}

% --- Le résumé --------------------------------------------------------------
%
% Il ouvre la lecture modernisée, sous le titre « Résumé », et il est composé
% en petit. Ce n'est pas un ornement : c'est le seul passage du document qui
% ne soit pas des mathématiques, et un lecteur doit pouvoir le reconnaître
% avant de l'avoir lu — puis le sauter s'il connaît déjà le sujet, ou s'y
% arrêter s'il ne le connaît pas. Le filet qui le suit marque la reprise du
% corps.

\newenvironment{resume}
  {\small\setlength{\parskip}{0.45em}}
  {\par\medskip\hrule\medskip}

% --- Mots-clés ---------------------------------------------------------------
%
% En anglais, en clôture du résumé de la lecture modernisée : le vocabulaire
% moderne sous lequel chercher ce que les feuillets construisent. Le manifeste
% du site (`npm run manifest`) les extrait pour étiqueter la cote entière —
% cette ligne est la source unique des tags, il n'y a pas de fichier à côté.
\newcommand{\keywords}[1]{\par\smallskip\noindent
  {\footnotesize\textsc{Keywords} --- \textit{#1}}\par}

% --- L'appareil critique ----------------------------------------------------

% Le numéro de feuillet, en marge. C'est l'ancre qui permet de régler un
% désaccord sur une formule en regardant *un* feuillet plutôt qu'une chemise
% de six cents. Le site s'en sert aussi pour tourner les pages du fac-similé
% quand on fait défiler la transcription.
\newcommand{\leaf}[1]{%
  \marginnote{\footnotesize\color{gray!70}\textsf{#1}}%
  \label{leaf:#1}\ignorespaces}

% Illisible : on ne devine pas. Un mot inventé qui se lit comme les autres est
% le pire résultat possible.
\newcommand{\ill}{\textcolor{red!60!black}{[\,\ldots\,]}}

% Lecture douteuse : le mot est donné, et signalé comme incertain.
\newcommand{\uncertain}[1]{\uline{#1}}

% Ajout de l'éditeur — une lettre manquante, un mot sous-entendu.
\newcommand{\add}[1]{\textcolor{blue!50!black}{[#1]}}

% Biffé par Grothendieck lui-même. Ce qu'il a rayé fait partie du document :
% une rature dit souvent où la pensée a changé de direction.
\newcommand{\struck}[1]{\sout{#1}}

% Note du transcripteur, sur l'état matériel du feuillet.
\newcommand{\note}[1]{\par\smallskip{\small\color{gray!60!black}\textit{[#1]}}\par\smallskip}

% Note marginale de Grothendieck, distincte du corps du texte.
\newcommand{\marginal}[1]{\par\smallskip\noindent\hspace{1em}%
  {\small\color{gray!60!black}#1}\par\smallskip}

% --- Titre ------------------------------------------------------------------

\AtBeginDocument{%
  \begin{center}
    {\large\bfseries Fonds Alexandre Grothendieck --- cote n\textsuperscript{o}~\grothFolder}\\[2pt]
    {\itshape\grothFoldertitle}\\[4pt]
    {\small feuillets \grothFirst--\grothLast\ (lot \grothBatch)}\\[2pt]
    {\footnotesize\color{gray!60!black}Transcription établie d'après le fac-similé numérisé
     par l'Université de Montpellier.\\
     Les lectures incertaines sont soulignées, les illisibles notées [\,\ldots\,],
     les ajouts entre crochets.}
  \end{center}
  \vspace{1em}}

\endinput


\folder{18}
\batch{2}
\pages{21}{40}
\dating{[à partir de 1969-vers 1972]}
\watermark{Édition de démonstration}
\foldertitle{Motifs et théorie de Hodge : notes manuscrites (s.d.)}

\begin{document}

\section{Généralisation de la formule fondamentale au cas d'un morphisme propre quelconque}

\note{Ces feuillets n'ont pas de couverture : le titre reprend mot pour mot la
première phrase de la page 22, qui n'est pas un titre de sa main. Ce qu'il
généralise est la formule fondamentale établie aux pages 16 à 20 du lot
précédent ; la page 21 est un feuillet étranger et l'argument ne s'interrompt
donc qu'en apparence.}

\page{22}

On généralise la formule fondamentale --- au cas d'un morphisme \emph{propre}
\emph{quelconque}
\[
  f : X \longrightarrow Y
\]
\ill{} \uncertain{variétés} \struck{\uncertain{algébriques}} holomorphes,
purement de codimension relative $p$ :

\begin{tikzcd}[column sep=small, row sep=small, nodes={font=\scriptsize}]
H^{2i}(X, \mathbb{C}) \arrow[r, "f^{\mathrm{top}}_{*}"] \arrow[d, "\wr"'] & H^{2i-2p}(Y, \mathbb{C}) \arrow[d, "\wr"] \\
H^{2i}(X, \Omega^{\bullet}_{X}) \arrow[r, "\frac{1}{(2i\pi)^{p}} f^{\mathrm{an}}_{*}"] & H^{2i-2p}(Y, \Omega^{\bullet}_{Y})
\end{tikzcd}

\note{Ce carré est encadré. Les exposants des $H$ sont repris à l'encre par
dessus une première écriture et ne se lisent pas avec certitude : à gauche
$2i$, à droite un exposant plus long où l'on croit lire $2i-2p$. Le décalage de
degré est celui qu'impose l'image directe d'un morphisme propre de codimension
relative $p$, mais la page ne le dit pas.}

D'où généralisation \ill{} \uncertain{analogue} pour des variétés
\uncertain{algébriq.} sur $\mathbb{C}$
\note{la phrase s'achève sur des points de suspension.}

\page{24}

$\mathbb{C}$, on a un objet canonique de $\xi$ \ill{} de $T_{\mathbb{C}}(\mathbb{Z})$
\note{la lettre entre « de » et $T_{\mathbb{C}}$ est surchargée et ne se lit
pas.}

\uncertain{[$T_{\mathbb{C}}$ torseur fondamental associé à} $\mathbb{C}$].

\[
  H^{i}(X, \mathbb{Q}(0)) \longrightarrow H^{i}(X_{\mathrm{top}}, \mathbb{Z})
\]

L'hom.\ $H^{i}(X'', \mathbb{Q}(0)) \to H^{i}(Y'', \mathbb{Q}(0))$
\uncertain{défini} \add{par une correspondance, à l'aide de} \ill{} par
$\xi \in a^{n}(X \times Y)$ est celui \uncertain{défini} par la classe
$P_{\mathrm{top}}(\xi) \left[ = \left(\frac{1}{2i\pi}\right)^{n} P_{\mathrm{alg}}(\xi) \right]$ :
\[
  \boxed{\ \xi(y)(x) = \mathrm{pr}^{\mathrm{top}}_{2*}\left( \mathrm{pr}^{*}_{1}(x)\, P_{\mathrm{top}}(\xi) \right)\ }
\]

Considérons l'\uncertain{isomorphisme} can\add{onique} (pour $X$, $i$ donnés)
\[
  \xi_{\mathbb{C}}\left(H^{i}(X, \mathbb{Q}(0))\right) = \xi\left(H^{i}(X, \mathbb{Q}(0))\right) \otimes_{\mathbb{Z}} \mathbb{C} \ \xrightarrow{\ \sim\ }\ H^{*i}(X, \Omega^{*}_{X}) = \xi_{\mathbb{C}, \infty}\left(H^{i}(X, \mathbb{Q}(0))\right)
\]
\note{sous le membre de gauche il écrit $H^{i}(X_{\mathrm{top}}, \mathbb{C})$
avec un signe $=$, et au-dessus de la flèche « car \ill{} Hodge-\uncertain{coh.} ».}

Cet isom.\ \struck{\ill{}} \uncertain{est} fonctoriel \ill{}
\struck{$H^{i}(X, \mathbb{Q}(0))$}.

En effet, \struck{\ill{}} \struck{Pour le rendre canonique, il faut} \ill{}
\[
  x \longmapsto \mathrm{pr}^{\mathrm{top}}_{2*}\left( \mathrm{pr}^{*}_{1}(x)\, P_{\mathrm{top}}(\xi) \right) = \mathrm{pr}^{\mathrm{top}}_{2*}\left( \mathrm{pr}^{\mathrm{an}}_{1}(x) \left(\tfrac{1}{2i\pi}\right)^{n} P_{\mathrm{alg}}(\xi) \right)
\]
\[
  = \left(\tfrac{1}{2i\pi}\right)^{-n} \mathrm{pr}^{\mathrm{alg}}_{2*}\left( \mathrm{pr}^{*}_{1}(x) \left(\tfrac{1}{2i\pi}\right)^{n} P_{\mathrm{alg}}(\xi) \right)
\]
\note{les exposants $-n$ et $n$ des deux facteurs $1/2i\pi$ sont écrits ainsi,
et le premier $\mathrm{pr}_{1}$ porte l'exposant « an » là où les autres portent
une étoile.}

\begin{tikzcd}[column sep=small, row sep=small, nodes={font=\scriptsize}]
H^{i}(X, \mathbb{C}) \arrow[r, "\xi_{\mathbb{C}}(\xi)"] \arrow[d, "\mathrm{can}"'] & H^{i'}(Y, \mathbb{C}) \arrow[d, "\mathrm{can}"] \\
H^{p}(X, \Omega^{\bullet}_{X}) \arrow[r, "\xi_{\mathbb{C},\infty}(\xi)"] & H^{i'}(Y, \Omega^{\bullet}_{Y})
\end{tikzcd}

\note{sous la flèche du haut il écrit « défini par
$P_{\mathrm{top}}(\xi) = (\frac{1}{2i\pi})^{n} P_{\mathrm{alg}}(\xi)$ »,
sous celle du bas « défini par
$P_{\mathrm{alg}}(\xi)$, $= (\frac{1}{2i\pi})^{n} \xi(\xi)$ » ; deux flèches
ondulées portent les deux lois déjà écrites plus haut, la seconde étant
$x \mapsto \mathrm{pr}^{\mathrm{alg}}_{2*}(\mathrm{pr}_{1}^{*}(x) P_{\mathrm{alg}}(\xi))$ avec « alg » souligné.}

\marginal{Cette \uncertain{vérification} \ill{} \uncertain{purement} analytique
$\longrightarrow$ topologique (avec \uncertain{chgt} de \uncertain{coord.}
\uncertain{analytiques})}

Ouf !

\page{26}

Soit $K$ un corps top.\ isomorphe à $\mathbb{C}$ [isom.\ non choisi !]

Alors on a le foncteur \add{mult.} $\xi = \xi_{K}$ sur la catégorie des motifs
sur $K$, à valeurs \ill{} les $\mathbb{Z}$-modules, et \ill{} isom.
\[
  \xi \otimes_{\mathbb{Z}} \mathbb{Z}_{\ell} \ \xrightarrow{\ \alpha_{1}\ }\ \xi_{\ell}
  \qquad
  \xi \otimes_{\mathbb{Z}} K \ \xrightarrow{\ \alpha_{2}\ }\ \xi_{K}
\]

Considérons
\[
  \xi\left(H^{2}(\mathbb{P}^{1}, \mathbb{Z}(0))\right) = \xi\left(\mathbb{Z}(-1)\right)
\]
c'est isomorphe à $\mathbb{Z}$, l'isom.\ n'étant \uncertain{pas} \ill{}
\add{ni déterminé les gén.} Le choix d'un tel isom.\ \uncertain{revient} à
celui d'une \emph{orientation} sur $\mathbb{P}^{1}_{K}$
\note{« orientation » est souligné, non biffé.}

[puisqu'en \ill{} \ill{} naturelle une variété \struck{algébrique} réelle,
\uncertain{iso.}, isomorphe à $S^{2}$), on a \ill{} $=$ \ill{}, au choix de
l'une des \struck{deux structures} \uncertain{orientées} sur
$\mathbb{P}^{1}_{K}$, provenant du \struck{choix d'un} choix d'un isom.\
$K \simeq \mathbb{C}$, i.e.\ le choix d'une \uncertain{des} \ill{}
\uncertain{racines} de $-1$ dans $K$. Comme
$\pi \in \mathbb{R} \subset \mathbb{R}$ \struck{$\mathbb{R}$} est bien
déterminé comme

\page{28}

\uncertain{élément} de $K$, on voit que le choix ainsi fait permet d'écrire
\[
  \alpha_{\infty}(u) = \frac{1}{2\sqrt{-1}\,\pi}\, v
\]
\ill{} $v$ est le générateur canonique de $H^{2}(\mathbb{P}^{1}, \Omega^{1})$.

Pour les choix \struck{différents} \uncertain{différents} de topologie sur $K$
\ill{} de $K$ \ill{}, on trouve ainsi \ill{} les éléments
\uncertain{primitifs} de $H^{u}(\mathbb{P}^{1}, \Omega^{u})$.
\note{les deux exposants sont écrits d'un même trait bref que rien ne permet de
lire.}

Si $K$ est un corps \struck{top.} isomorphe à $\mathbb{R}$ [l'isom.\ étant
d'ailleurs d'ailleurs \emph{unique}], on définit encore, par choix de la
\struck{compactification} clôture alg.\ \uncertain{algébrique} de $\mathbb{R}$,
un foncteur $\xi = \xi_{K} = \xi_{K}$ à valeurs dans les $\mathbb{Z}$-modules
\add{mais avec op.\ de $\mathbb{Z}/2\mathbb{Z} = \mathrm{Gal}(\overline{\mathbb{R}}/\mathbb{R})$},
\struck{et \ill{} des isom.} un dual sur $\mathbb{Z}$, les
$\xi_{\otimes}$-\uncertain{réseaux} \add{$\xi(M)$} pour \ill{},
\uncertain{obtenons} dans
\[
  \prod_{\ell} M_{\ell} = M_{\ell} \times M_{\infty \mathbb{C}}
\]
\add{valeurs : $\mathbb{R} \hookrightarrow \mathbb{C}$, $\mathbb{Z}/2\mathbb{Z}$}
\ill{} \ill{}
\note{La page s'arrête sur une ligne de deux ou trois mots que rien ne permet
de lire. Le produit $\prod_{\ell} M_{\ell} \times M_{\infty\mathbb{C}}$ est le
premier état, sur ces feuillets, du module adélique que le lot 3 écrira
$M^{\mathbb{A}_{\mathbb{C}}}$.}

\section{Corps de base topologique $\simeq \mathbb{R}$ ou $\mathbb{C}$ (foncteurs de Hodge et De Rham dans ce cas)}

\note{Le titre est de lui : ces quatre lignes occupent seules le haut à droite
du feuillet 30, par ailleurs vierge. Ce feuillet sert de couverture à ce qui
suit et n'est pas transcrit comme page, d'où le saut de 28 à 31. « typ. » y
abrège « topologique », et « De » est biffé puis réécrit dans « DeRham ».}

\page{31}

\struck{Corps} Corps de base topologique donné : $\mathbb{R}$ ou $\mathbb{C}$
\note{cette ligne, soulignée, est en tête du feuillet ; elle reprend le titre
de la couverture.}

$K$ corps topologique donné $\simeq$ \ill{} $\mathbb{C}$ (il y a donc
\emph{deux} isom.\ avec $\mathbb{C}$, mais on \ill{} on en choisit
\uncertain{un}) $j$. \ill{} cas $\mathbb{R}$.

$M$ motif sur $K$. Alors \struck{$D$} \struck{$DR(M)$} \add{De} \add{De Rham}
pour \ill{} $X$ \uncertain{lisse projectif} sur $K$, $X(K)$ est une \ill{}
variété topologique \ill{} ($K$-analytique compacte), d'où d'où
$H^{*}(X, \mathbb{Q})$, et on définit ainsi \struck{le foncteur} le
\emph{Betti} \emph{Betti} $\mathrm{Betti}(M)$ sur les motifs sur $K$.
$\mathrm{Betti}(\mathrm{Betti}(M) \otimes_{\mathbb{Q}} K$ \struck{$\simeq$}
correspond à $H^{*}(X, K)$. Par l'argument \uncertain{par} \uncertain{habituel},
ce dernier est isomorphe canon.\ \ill{} la coh.\ coh.\ de De R., d'où
\[
  \mathrm{Betti}(M)_{K} \simeq DR\struck{(\subseteq DR(M)} \qquad \text{($\otimes$-isom.\ can.)}
\]
\note{la parenthèse du membre de droite est biffée en cours d'écriture et
reprise en $DR(M)$ ; c'est cette dernière forme qui est utilisée plus bas.}

Ceci met sur \add{(à $K$-modules)} $DR(M)$ \struck{$(M)$} une structure
\emph{réelle}, dont le sous-\ill{} \uncertain{réel} \uncertain{réel} est
$\mathrm{Betti}(M)_{\mathbb{R}}$ ; sauf à l'\uncertain{isomorphisme}. La
\uncertain{donnée} de la filtration canonique de $DR(M)$, et $\sigma$,
permettent \struck{\ill{}} de \struck{façon} unique de \uncertain{splitter}
cette filtration \struck{\uncertain{-cation}}, de \uncertain{façon} à
déterminer \ill{} $DR(M)$ une \struck{bi} \emph{bigraduation} compatible avec
les degrés \uncertain{totaux}, et \ill{} \uncertain{redonnant} la filtration
\ill{} à l'aide des \uncertain{premiers} degrés. Le choix d'un isom.\
\uncertain{$u$} $K \to \mathbb{C}$ met, \uncertain{pourtant}, pas
\uncertain{beaucoup} de \ill{}, une

\page{32}

\add{le naturel} structure \struck{\uncertain{complexe}} sur \struck{$I$}
$DR(M)$ ; mais \ill{} \uncertain{noter} que la bigraduation déduite
\uncertain{ainsi} par \uncertain{transport} de \uncertain{structure} de celle
de Hodge \uncertain{ne} \ill{} \emph{dépend pas de} \ill{} (donc, si on
\uncertain{remplace} \ill{} $u(z)$ \uncertain{par} $\overline{u(z)}$, elle
\uncertain{n'est} \emph{pas} \uncertain{topologie} \ill{} la
\uncertain{symétrique} !).

\uncertain{un troisième} \uncertain{$u(K)$} \ill{}
$DR(M_{\mathbb{C}}) \xrightarrow{\ \sigma\ } M_{\mathbb{C}} \xrightarrow{\ \sigma\ } DR(M)_{\mathbb{C}}$,
les structures \ill{}
\struck{\ill{} la filtration} composé avec $\sigma$, ou $\overline{\sigma}$,
\uncertain{introduisant} un \emph{$K$-isom.}
$DR(M_{\mathbb{C}}) \simeq DR(M), DR(M)$, conservant les
$\mathbb{Q}$-structures, et \struck{\uncertain{renversant}} \add{(à filtration)}
\uncertain{filtration}.

\marginal{Néanmoins, si on regarde, pour un motif donné $M$, le motif
$\otimes(M_{\mathbb{C}})$ ($\subseteq$ à l'autre), les \ill{} \uncertain{ne}
\ill{} \uncertain{sont} \ill{}}

\struck{Supp.} Soit maintenant \struck{$K$} $K_{0} = \mathbb{R}$, et
\uncertain{considérons} les motifs $M$ \struck{sur} $\mathbb{R}$. Choisissons
\ill{} \uncertain{une} clôture \uncertain{alg.} $K$ de $K_{0} = \mathbb{R}$.
\struck{Alors} Alors
\[
  DR(M)_{K} = DR(M_{K})
\]
est muni de \emph{deux \struck{fois} façons} d'une $\mathbb{R}$-structure :
\uncertain{la} \uncertain{première} \uncertain{déduite} de $DR(M)$ par
$\mathbb{R} \to K$, et \uncertain{puis} : la \uncertain{contre-structure}
\uncertain{donnée} plus haut. Soit $\sigma$ l'aut\add{omorphisme} de la
\uncertain{conjugaison} \uncertain{complexe} \uncertain{définie} plus haut,
(N.B.\ $\sigma$ \ill{} $\sigma$ ne \emph{respecte pas} la
\struck{\uncertain{bigraduation}} \add{filtration}), $\tau$ celle
\uncertain{déduite} par \ill{} $DR(M)_{K} \simeq DR(M_{K})$ (donc $\sigma$
\uncertain{respecte} la filtration $W$ \ill{} $W$) --- $\sigma$ et $\tau$ sont
\struck{\uncertain{semi-linéaires}}. \struck{Alors} Alors $\sigma$ et $\tau$
commutent, en d'autres termes \ill{}
\[
  DR(M_{K})^{\tau}\ \left(= H^{*}(X(K), \mathbb{R}) \text{ si } M = M^{*}(X)\right)
\]
est stable par $\tau$ : en \uncertain{effet}, $\tau$ n'est autre \ill{} par
l'aut\add{omorphisme} de \ill{} $H^{*}(X(K), \mathbb{R})$ défini par
l'aut\add{omorphisme} $\tau_{X}$ de \ill{} $X(K)$ ! Alors
$\sigma\tau = \tau\sigma = f_{\mathbb{R}}$ \ill{} l'aut\add{omorphisme} de
\struck{$H^{\bullet}$} $H^{*}(X(K), K) \simeq DR(\mathrm{Mot}^{*}(X))$

\page{33}

\uncertain{défini} par l'aut\add{omorphisme} \struck{$\tau_{X}$} $\tau_{X}$ de
$X(K)$). Ainsi, un motif sur $\mathbb{R}$ définit \ill{} :

\begin{enumerate}
  \item Une structure de \struck{Hodge} Hodge \add{$V$} $V$ sur $K$
  \item Un automorphisme \struck{\ill{}} \add{devant devant} \ill{}
        \add{$V_{K}$} \ill{} $K$ --- qui \uncertain{est} \ill{}, cet
        automorphisme étant étant « réel » i.e.\ \uncertain{qui} transforme
        $V_{\mathbb{R}}$ en $V_{\mathbb{R}}$, et \ill{} \emph{antilinéaire}
        \uncertain{et} i.e.\ transforme $V$ en $\overline{V}$
\end{enumerate}

\struck{\ill{}} Ces deux dernières étant \uncertain{tant}
\uncertain{assujetties} à la condition

\begin{quote}
($f_{\mathbb{R}}$) \quad $\sigma$ \ill{} \uncertain{échange} la filtration
naturelle, i.e.\ \uncertain{qu'il} échange échange \uncertain{que} la
filtration \uncertain{donnée} et l'\uncertain{opposée}
\end{quote}

En termes de « \uncertain{bonnes} \uncertain{structures} de \uncertain{Hodge} »,
on peut \struck{\ill{}} \add{exprimer} aussi b) en disant \uncertain{qu'il}
\uncertain{s'agit} \ill{} d'un \uncertain{isom.} de la structure de
\struck{Hodge} Hodge $V$ avec la structure de Hodge \struck{\ill{}}
« conjuguée » \add{« $-$ »} : c'est le \emph{Frobenius réel}.

Soit $G$ le groupe de \ill{} de Galois \uncertain{motivique} de $M$ sur
$\mathbb{R}$. Alors $f_{\mathbb{R}}$ correspond à \ill{} un élément
\[
  f_{\mathbb{R}} \in G(\mathbb{Q}) \ \struck{G(\mathbb{Q})}
\]
« Weil » \emph{Frobenius réel} \emph{réel}. C'est un élément d'ordre 2.

Son image dans $G/G^{0}$ \struck{$G/G^{0}$} engendre ce groupe. Cet élément
normalise $i : S_{\mathbb{R}}$ \struck{$\to$} \struck{$: S_{\mathbb{R}}$}
$\longrightarrow G_{\mathbb{R}}$, de façon précise
\[
  f_{\mathbb{R}}\, i(\lambda)\, f_{\mathbb{R}}^{-1} = i\left(\struck{\ill{}}\, i(\struck{\ill{}}\, \lambda')\right)
\]
\marginal{où $\lambda \mapsto \lambda'$ est la \uncertain{symétrie} canonique
de $S_{\mathbb{R}}$}
\note{la formule est écrite deux fois l'une sur l'autre ; seule la seconde
version, avec $\lambda'$, est cohérente avec la note marginale et avec la page
36.}

Donc sur $U_{\mathbb{R}}$ \ill{}
\[
  (*) \qquad f_{\mathbb{R}}\, i(\lambda)\, f_{\mathbb{R}}^{-1} = i(\lambda^{-1}) \qquad (\lambda \in U_{\mathbb{R}}\ \struck{\ill{}})
\]

Le foncteur $M \mapsto$ (\struck{struct.} structure de Hodge
\uncertain{associée}, avec $f_{\mathbb{R}}$) est \uncertain{pleinement} fidèle
\ill{} \ill{} \ill{} des \uncertain{conj.} de Hodge.

\page{34}

Les éléments de $DR(M)$ \struck{$(M)$} \uncertain{sont} les
\uncertain{éléments} de $DR(M)_{K} \simeq V \otimes_{\mathbb{Q}} K$ \ill{} $z$
de la forme $z = x + iy$ \quad ($x, y \in V_{\mathbb{R}}$) avec $x = fx$,
$y = -fy$
\add{(ce qu'on exprime $\tau z = z$ i.e.\ $\sigma\tau z = z$ i.e.\ $\sigma z = f(z)$)}
\marginal{$\mathbb{Q}$-structure de \ill{}}

\struck{Cela} \struck{$\pm$} Donc l'\uncertain{espace} $DR($
\struck{$DR(M)$} admet une \struck{\uncertain{façon}} \uncertain{naturelle}
\ill{} \struck{\ill{} structure} \ill{} \uncertain{peut} \uncertain{peut} le
\ill{} \ill{} en $x + iy$, $x, y \in V$, $x = fx$, $y = -fy$.
\struck{$-fy$.} Soit $W$ cet \uncertain{espace}, il mérite le
\struck{nom} de \emph{sous-\ill{} \uncertain{espace} de Betti} \add{tordu} de
$DR(M)$.

\begin{tikzcd}[column sep=small, row sep=small, nodes={font=\scriptsize}]
H^{*}(X(K), \mathbb{Q}) = V \arrow[r, hook] & DR(M)_{K} \simeq DR(M_{K}) \\
H^{*}(X(\mathbb{R}), \mathbb{Q}) = V \cap W = V^{f} \arrow[u, hook] & W \arrow[u, hook]
\end{tikzcd}

\note{Il écrit les deux inclusions verticales par des $\cup$ ; le membre de
gauche de la ligne du bas porte encore « $= W^{f}$ » suivi d'un
\struck{$\mathbb{C}$}.}

\marginal{Il est stable par $f$. Mais la formation \ill{} \uncertain{ne}
\ill{} \uncertain{commute-t-elle} \uncertain{comme} il faut \uncertain{pour}
$\otimes$ ? \emph{Oui}.}

On a ainsi un \struck{deuxième} deuxième \uncertain{foncteur} fibre \ill{} $M$
\add{\uncertain{un foncteur} \uncertain{Betti}} : \ill{} sur $\mathbb{Q}$,
\struck{$\mathbb{Q}$}, mais il est \uncertain{isomorphe} \ill{}
\uncertain{premier}, \struck{\ill{}} \ill{} \ill{} \ill{} \ill{}
\[
  V \simeq V^{+} + V^{-}, \qquad W = W^{+} + W^{-+} + W^{-},
\]
et \uncertain{utilisant} les isomorphismes
$V^{+} \simeq W^{+}$ \struck{$\simeq W^{+}$}, $V^{-} \simeq W^{-}$ canoniques.

Les isom.\ $V \to W$ \ill{} \add{$f$.} \ill{} par l'\uncertain{automorphisme}
de $DR(M_{K}) \simeq V_{K} \simeq W_{K} \cdot W_{K}$ qui est l'identité sur
$V^{+}_{K}$, et la multiplication \uncertain{par} \uncertain{$i$} sur
$V^{-}_{K}$ ; \add{Cependant, ce} Donc \uncertain{n'est} \emph{pas}
\struck{compatibl.} \uncertain{compatible} avec \ill{} \ill{}

\add{\ill{} Betti et $\Pi$ et \struck{Betti} un \ill{} \ill{}
\uncertain{isom.}, l'un est déduit \ill{} de l'\ill{} de l'autre en
\uncertain{tordant} par le cocycle \uncertain{défini} par $f_{\mathbb{R}}$}
\note{« Betti » est réécrit une seconde fois au-dessus de « déduit ».}

\page{35}

\note{Ce feuillet et les deux suivants sont à l'encre bleue et écrits au net ;
ils se lisent presque de bout en bout, à la différence des pages 24 à 28.}

L'élément $\alpha \in P(\mathbb{C})$ est \ill{} est \uncertain{relié} à
$f_{\infty}$ par
\[
  \overline{\alpha} = f_{\infty}\, \alpha
\]
donc \struck{$f_{\infty}$} à priori
\[
  f_{\infty} = 1 \iff \struck{\iff} \alpha \in P(\mathbb{R}) \ \text{ i.e. } \ \mathbb{R}(\alpha) = \mathbb{R}
\]

\struck{Comme on a la relation $f$ \ill{} $f_{\infty}$ et $f^{-1}_{\infty} = \mu$, cette condition implique $f_{1} = f_{2}$ \ill{} $f_{1} = f_{2} = \frac{1}{2}i$, i.e.\ \ill{} dont conj.\ de Hodge $f$ \ill{} $G^{0} \xrightarrow{\ \varepsilon\ } G_{m}$}
\note{tout ce passage, quatre lignes, est barré d'un long trait horizontal ;
ce qui suit le reprend sous une autre forme.}

Or on vérifie aussitôt tôt
\[
  \boxed{\ \varepsilon(f_{\infty})\, \varepsilon(f_{\infty}) = -1\ }
\]
ce qui nous montre qu'on qu'on --- \uncertain{types}
\[
  \boxed{\ f_{\infty} \neq 1 \neq 1 \ \text{ i.e. } \ \mathbb{R}(\alpha) = \mathbb{C}\ }
\]
(si $M$ \uncertain{contient} le motif de Tate)

Si $M$ \uncertain{contient} \ill{} le \struck{le} motif de Tate, on voit que
$f_{\infty} = 1 \Rightarrow f_{1} = f_{2}$ \struck{$1 = f_{2}$} i.e.\
$f_{1} = f_{2} = \frac{1}{2}i$, donc en \uncertain{vertu} de la
\uncertain{conj.}\ de \struck{Hodge} Hodge, il faut que \ill{} ou bien bien
$G^{\delta} = \{e\}$ \struck{\ill{}} $G^{0} = \mathrm{Im}\, f_{1}$, donc ou
bien ici \struck{$G$} ici $G = G$ \uncertain{isogène} \ill{}
$\mathbb{Z}/2\mathbb{Z}$, donc \ill{} \struck{$G = \{e\}$}
\struck{$= \{e\}$} \ill{}

[N.B.\ si $G$ si $G$ est \uncertain{discret}, on sait que \ill{} \ill{} \ill{}
$G = 1$ \ill{} \struck{donnera on $f_{\infty}$} bien ici
$G = G_{m}/\mu_{n}$, $n \geqslant 1$, $f_{\infty} \notin G^{0}(\mathbb{Q})$],
\ill{} $\mathbb{Q}(-n)$ ; on \uncertain{trouve}

$M$ est \struck{\ill{}} engendré engendré par \ill{} \ill{}
\uncertain{bien}. Donc on trouve \ill{} $=$ \ill{} que \ill{}
\begin{align}
  \left(f_{\infty} = 1\right)
    &\iff \left(\mathbb{R}(\alpha) = \mathbb{R} \ \struck{\mathbb{R}(\alpha) = \mathbb{R}}\right) \notag \\
    &\iff M \text{ est engendré par } \mathbb{Q}(-2n), \text{ pour quelques } n \geqslant 0 \notag
\end{align}
\[
  \Updownarrow
\]
\[
  G \simeq G_{m} \quad \text{ou} \quad G \simeq \{e\}
\]

\marginal{N.B.\ $f_{\infty}$ ne peut être \emph{autre} que si
$G^{0} = \mathrm{Im}\, \varepsilon i$ \struck{$= \mathrm{Im}\, \varepsilon i$},
ce qui implique que $M$ est \struck{\ill{}} \add{contenu dans \ill{}} \ill{}
engendré par $\mathbb{Q}$ \uncertain{tordu} et par \ill{} et par
$\mathbb{Q}(2)$. Je \uncertain{reviens} \ill{} dire que $G$ est ce $G$ est
comm.}
\marginal{\uncertain{uniquement} \uncertain{conj.}\ de Hodge}

\page{36}

\emph{Motifs sur $\mathbb{R}$}

$M$ catégorie \add{$\otimes$} des motifs sur $\mathbb{R}$, $G = \mathbb{R}$,
$G = \mathrm{Aut}_{\otimes}(T_{B})$ groupe de Galois associé.
\note{Il souligne $\mathrm{Aut}_{\otimes}$, comme plus bas $\mathrm{Isom}_{\otimes}$ et $\mathrm{Norm}$ : c'est sa façon de marquer le foncteur, et non une lecture douteuse. Le soulignement n'est pas reproduit, la transcription le réservant à l'apparat.}

Notons qu'on a un $\otimes$-automorphisme \struck{automorphisme} d'ordre 2 de
$T_{B}$, qui sur la catégorie $\mathcal{V}_{\mathbb{R}}$
\struck{$\mathcal{V}_{\mathbb{R}}$} des var.\ proj.\ lisses $X$ sur
$\mathbb{R}$ revient à l'automorphisme de \struck{des}
$H^{*}(X(\mathbb{C}), \mathbb{Q})$ défini par l'autom.\
$x \mapsto \overline{x}$ de $X$ \struck{$X(\mathbb{C})$}. D'où
\[
  \boxed{\ f_{\mathbb{R}\infty} \in G(\mathbb{Q})\ \struck{G(\mathbb{Q})}\ } \qquad (f_{\infty}^{2} = 1)
\]
(« frobenius réel »). On \struck{retrouve} retrouve
\struck{\uncertain{le sous-espace}}
\struck{$H^{*}_{DR}(X)$ de $H^{*}(X, \mathbb{C})$} $= +$ si $+$ si
$\sigma_{M}$ est la conjugaison complexe de $M \otimes_{\mathbb{Q}} \mathbb{C}$
[i.e.\ \struck{i.e.} $M \in \mathrm{ob}\,\mathrm{Rep}(G)$], alors la
conjugaison $\tau_{M}$ associée \struck{associée} à la structure réelle
$T_{DR}(M) \subset M \otimes_{\mathbb{Q}} \mathbb{C}$ n'est autre
\struck{autre} que
\[
  \tau_{M} = (f_{\infty})_{M}\ \struck{(f_{\infty})_{M}}\ \sigma_{M} = \sigma_{M}\, (f_{\infty})_{M}
\]

Donc on retrouve qu'on a $f_{\infty}$ \struck{$f_{\infty}$} la structure mixte
sur $M \otimes_{\mathbb{R}} \mathbb{C}$, donc aussi \ill{} le torseur
\[
  P = \mathrm{Isom}_{\otimes}\ \struck{\mathrm{Isom}_{\otimes}}\ \left(T_{B} \otimes_{\mathbb{Q}} \mathbb{R},\ T_{DR}\right)
\]
sur $G_{\mathbb{R}}$ : c'est celui défini \struck{défini} par le 1-cocycle
\struck{\ill{}} \uncertain{normalisé} sur $\Gamma = \mathbb{Z}/2\mathbb{Z}$, à
valeurs \struck{dans} \struck{dans} $G(\mathbb{C})$, dont la valeur en
l'élément $\neq e$ est $f_{\infty}$. \struck{La} La forme tordue correspondante
de $G$ s'obtient en \struck{prenant} prenant sur $G_{\mathbb{C}}$ la \ill{} de
\uncertain{descente} associée : \struck{\ill{}} \struck{\ill{}}
$\mathrm{int}(f_{\infty}) \times_{\mathbb{R}} \sigma_{\mathbb{C}}$.
\struck{On a dit}

\struck{donc, pour} La relation entre \struck{$f_{\infty}$} $f_{\infty}$ et
\[
  j : S_{\mathbb{R}}\ \struck{\to S_{\mathbb{R}}}\ \longrightarrow G_{\mathbb{R}}
\]
est que
\[
  \boxed{\ f_{\infty}\, j(\lambda)\, f_{\infty}^{-1} = \struck{j(\ill{})} = j(\uncertain{\lambda'})\ }
\]
\marginal{i.e.\ $\lambda \mapsto \lambda'$ est l'automorphisme can.\ de
$S_{\mathbb{R}}$ (donnant $\lambda \mapsto \overline{\lambda}$ sur
$S_{\mathbb{R}}(\mathbb{R}) \simeq \mathbb{C}^{*}$)}

\page{37}

D'autre part, on aura ici \struck{ici}
\[
  \beta \in Q(\mathbb{R}) \subset P'(\ \struck{)} \subset P'(\mathbb{R})
\]
donc
\[
  \mathbb{R}(\gamma) = \mathbb{R}(\beta, \gamma) = \mathbb{R}(\beta, \alpha) = \mathbb{R}(\beta, \alpha) = \mathbb{R}(\alpha)
\]
et par suite
\[
  \mathbb{R}(\gamma) = \mathbb{R}(\alpha)
\]
qui est donc $= \mathbb{C}$ \struck{sauf} \struck{sauf} dans le cas
exceptionnel précisé plus haut.

\emph{Remarque} Soit $M$ un motif sur $\mathbb{R}$, \struck{$\mathbb{R}$,}
correspondant à une \struck{\ill{}} \struck{\ill{}} $T_{B}(M)$. Alors on
\struck{on} ne peut avoir $f_{\infty M} = \mathrm{id}_{M}$ que si $M$ est une
somme \uncertain{directe} de puissances tensorielles de $\mathbb{Q}(2)$.

\note{Ce qui suit, jusqu'au bas du feuillet, est encadré puis annulé par trois
longs traits diagonaux. Le calcul se lit néanmoins et il est transcrit.}

La relation $f_{\infty}\, j(\lambda)\, f_{\infty}^{-1} = j(\lambda')$
\struck{$j(\lambda')$} indique, faisant
$\lambda = C \in S_{\mathbb{R}}$ \struck{$C \in S(\mathbb{R}$}
\[
  f_{\infty}\, C\, f_{\infty}^{-1} = C' = C^{-1} = \eta\, C \qquad (\eta = i(-1))
\]
d'où
\[
  f_{\infty} \in \mathrm{Norm}_{G_{\mathbb{R}}}\left(\mathrm{im}_{G_{\mathbb{R}}}(C, C^{-1})\right)
\]
\struck{$K' = \mathrm{Norm}_{G_{\mathbb{R}}}(C, C^{-1}) \supset \mathrm{Norm}_{G_{\mathbb{R}}}(C) = K$}
\ill{} $[= \mathrm{Ker}\,\varepsilon\,/\,\dot{i}(\mu_{2})]$, et \ill{} les
images $\dot{f}_{\infty}$, $\dot{C}$ de $f_{\infty}$, $C$,

\uncertain{Considérons} $G/i(\mu_{\ill{}}) = \dot{G}$

donc
\[
  \dot{f}_{\infty} \in \mathrm{Norm}\ \struck{\mathrm{Norm}}_{\dot{G}_{\mathbb{R}}}\left(\dot{C}, \dot{C}^{-1}\right) \overset{\mathrm{df}}{=} K'
\]
\[
  \dot{K} = \mathrm{Norm}\ \struck{\mathrm{Norm}}_{\dot{G}_{\mathbb{R}}}\left(\dot{C}\right) \subset K'
\]

\page{38}

\struck{\ill{}}

Soit $K$ un corps topologique \struck{-gique} isomorphe à $\mathbb{C}$, $X$ un
schéma \uncertain{algébrique} défini \struck{sur} sur $K$. Alors
$X_{\mathrm{top}}$ est défini, \add{si $X$ est non sing.} et aussi
$H^{*}(X_{\mathrm{top}}, K)$. \struck{$(\ )$} D'autre part, \add{la cohomologie}
de Hodge est définie, \ill{}, soit $H^{*}(X)$ [indépendant de la structure
topologique \struck{\ill{}} de $K$]. Alors le choix d'un isom.\ top.\
$u : K \to$ \struck{$\to$} $\mathbb{C}$ définit un isomorphisme
\[
  H^{*}(X_{\mathrm{top}}, K) \ \xrightarrow{\ \sim\ }\ \struck{(K)} \ \xrightarrow{\ \sim\ }\ H^{*}(X)
\]
et on vérifie aussitôt \uncertain{que} \uncertain{cet} isomorphisme ne dépend
pas du choix de \struck{de} $u$. [En d'autres termes, si $K = \mathbb{C}$,
alors les \struck{de} deux homomorphismes

\begin{tikzcd}[column sep=small, row sep=small, nodes={font=\scriptsize}]
H^{*}(X_{\mathrm{top}}, \mathbb{C}) \arrow[r] \arrow[rd] & H^{*}(X) \arrow[d] \\
{} & H^{*}(X^{c})
\end{tikzcd}

\note{la flèche verticale porte à droite « (antilinéaire \struck{linéaire}) »
et à gauche la mention « non commut.\ ».}

(où $c : \mathbb{C} \to \mathbb{C}$ est la conjugaison complexe) sont
\struck{\uncertain{conjugués}} \uncertain{déduits} l'un de l'autre
\uncertain{par} \uncertain{extension} de $c$ aux \struck{aux} coefficients de
$H^{*}(X_{\mathrm{top}}, \mathbb{C})$]

Cela nous permet, notant \struck{\ill{}} $u : \mathbb{R} \to \mathbb{C}$, de
définir \struck{$H^{*}(X_{\mathrm{to}}$}
$H^{*}(X_{\mathrm{top}}, \mathbb{R}) \hookrightarrow H^{*}(X)$ (même si $X$ \ill{} \uncertain{définie})
\add{un homom.\ injectif contenu \ill{}}

\struck{Supposons} $X$ projective \struck{et lisse} \ill{} non singulière,
alors le choix d'un \struck{isom.} \struck{isom.} top.\
$u : K \xrightarrow{\ \sim\ } \mathbb{C}$ définit une décomposition \ill{}
\[
  \varphi : \coprod_{p,q} H^{q}(X, \struck{H^{q}(X,}\ \Omega^{p}_{X}) \ \xrightarrow{\ \sim\ }\ H^{*}(X)
\]
compatible avec la \struck{filt.} filtration de $H^{*}(X)$. On vérifie encore
que $\varphi$ \struck{ou $\varphi$} est indépendant du choix de $u$.

\marginal{\struck{\ill{}} \struck{$\mathbb{R}$} \uncertain{notant}
$H^{*}(X_{?}, \mathbb{R})$ \struck{\ill{} de \ill{} $K$)} anti\ill{}
$H^{*}(X)^{-} = \sqrt{-1}\, H^{*}(X)^{+}$}

\page{39}

Enfin, supposons maintenant que $K$ soit isomorphe \struck{à} $\mathbb{R}$
[on peut supposer, ou $K = \mathbb{R}$, ce $\mathbb{R}$ n'ayant pas
d'automorphismes non triviaux] \struck{]} soit $K'$ une clôture algébrique de
$K$, $\pi = \mathbb{Z}/2\struck{;} = \mathbb{Z}/2\mathbb{Z}$ son groupe de
Galois, $X$ lisse [\struck{\uncertain{projectif}}] sur $K$ \struck{$K$}. Alors
on peut expliquer \uncertain{ce} \uncertain{qui} précède :
\struck{$H^{*}(X_{\mathrm{top}}, \mathbb{R}) \xrightarrow{\ \sim\ } H^{*}(X) \otimes_{K} K'$}, on trouve
\note{cette ligne est écrite puis entièrement biffée, une première parenthèse « $, \mathbb{R})$ » étant déjà raturée en cours d'écriture ; ce qui suit la remplace.}
\[
  \left[\ \varphi_{K'} : \coprod_{p,q} H^{q}(X, \struck{\Omega\,}(X, \Omega^{p}_{X}) \otimes_{K} K' \ \xrightarrow{\ \sim\ }\ H^{*}(X \struck{\ }) \otimes_{K} K'\ \right]
\]
et d'après ce qui \ill{} été vu plus haut, \struck{ces} ces isomorphismes
\struck{\ill{}} \emph{commutent à $\pi$}, donc proviennent d'un isomorphisme
\ill{}
\[
  \varphi : \coprod_{p,q} H^{q}(\struck{;}\ H^{q}(X, \Omega^{p}_{X}) \ \xrightarrow{\ \sim\ }\ H^{*}(X)
\]

En résumé, chaque fois que $X$ est projectif non singulier sur un corps
\struck{top.} top.\ loc.\ compact \struck{compact} archimédien, on a une
décomposition $\varphi$ \uncertain{comme} dessus de $H^{*}(X)$, \ill{},
compatible avec la filtration. Si on \uncertain{pose} \struck{pose}
\[
  H^{p,q}(X) = \varphi\left(H^{q}(X, \Omega^{p}_{X})\right) \subset F^{p} H^{q}(X)
\]

Si $K$ n'est pas \struck{pas} topologique, alors pour \uncertain{tout}
\uncertain{choix} d'une \struck{une} topologie \uncertain{localement}
archimédienne $T$ sur $K$, on \uncertain{déduit} \ill{}

\marginal{\ill{} $\mathbb{Z}/2\mathbb{Z}$) \ill{} opérant via son opération sur
$X_{\mathbb{C}\,\mathrm{top}}$ \uncertain{compact} \emph{et} sur $\mathbb{C}$
dans $\otimes_{\mathbb{R}} \mathbb{C}$, $\simeq H^{*}(X_{\mathbb{C}\,\mathrm{top}}, \mathbb{R})^{\pi} \simeq H^{*}(X_{\mathrm{top}} = X_{\mathbb{C}\,\mathrm{top}}/\pi_{0}\,;\, \mathbb{R})$, \struck{$= H^{*}(\xi, \mathbb{C}\,\mathrm{top}, \mathbb{R})^{-}$}}

\page{40}

\[
  H^{*}(X)_{T} = \struck{H^{*}(X} = H^{*}(X) \otimes_{K} \widehat{K}_{T}.
\]
une décomposition :
\[
  \boxed{\ \varphi_{T} : \coprod_{p,q} H^{q}(X, \struck{\Omega\,}(X, \Omega^{p}_{X})_{T} \ \longrightarrow\ H^{*}(X)_{T}\ }
\]

Cela permet de \struck{défin.} définir
\[
  H^{p,q}_{T}(X) = \varphi_{T} \struck{=} \varphi_{T}\left(H^{q}(X, \Omega^{p}_{X})_{T}\right) \cap H^{*}(X) \ \subset\ F^{p}\!\left(H^{q}(X)\right)
\]
qui est un sous-\emph{$K$-espace} \struck{espace} \emph{vectoriel de
$H^{q}(X)$}, défini canoniquement \struck{par} par la donnée de $T$. Nous
pouvons de \struck{$=$}
\[
  H^{p,q}_{T}(X, \mathbb{R}) = \varphi_{T} \struck{=} \varphi_{T}\left(H^{q}(X, \Omega^{p}_{X})_{T}\right) \cap H^{*}(X_{T}, \mathbb{R}) \cap H^{*}(X)
\]
\struck{c'est \ill{} \ill{} un sous-espace vectoriel de $H^{q}(X)$}
c'est \ill{} un \ill{} $\widehat{K}_{T}$ : \struck{$\widehat{K}_{T} \simeq \mathbb{C}$}
si $\widehat{K}_{T} \simeq \mathbb{R}$, \ill{} \ill{} (c'est un \ill{}
\struck{$\mathbb{R}_{T}$} \uncertain{$\mathbb{R}_{T}$}-espace vectoriel), si
\ill{} si
\[
  \mathbb{R}_{T} = \struck{\mathbb{R}_{T}} = \mathbb{R} \cap K, \qquad \text{donc } \widehat{K}_{T} \simeq \mathbb{C}
\]
[sous-\emph{corps} \struck{\ill{}} \emph{réel} pour $T$ \struck{$r$ $T$}] ; à
\uncertain{fortiori}, \ill{} \ill{} \ill{} on \ill{} obtient
\uncertain{ainsi} \uncertain{ainsi} \struck{\ill{}} \ill{}, \ill{}
\[
  K_{r} \simeq \bigcap_{T} \mathbb{R}_{T} \ \struck{= \mathbb{R}_{T}} \ \simeq\ \text{plus grand sous-corps tot.\ réel de } K
\]
[N.B.\ $K_{r}$ est une ext.\ \struck{algéb.} algébrique \ill{}
\uncertain{corps} \uncertain{premier} $\mathbb{Q}$.]

\note{Le feuillet s'arrête là, aux deux tiers de la page. Rien ne se poursuit
au feuillet 41, qui est une couverture de sa main ouvrant un tout autre
développement (lot suivant).}

\end{document}
